\section{Methodology}\label{sec:Model}

Let $y_t$ be a variable that measures a teacher's subjective evaluation of a student in time period $t$, $SES$ a variable that measures the socioeconomic status of the student, and $s_t$ a variable that reflects student ability in time period $t$. The long regression equation equals,
\begin{align} \label{eqn:long}
    y_{t}= \beta SES + \gamma s_{t}+e_{t}.
\end{align}

Throughout the paper we suppress the subscript for student $i$ and the intercept in regression equations. Scholars and policy makers are generally interested in the parameter $\beta$. That is, they are interested in whether subjective teacher evaluations differ systematically by SES, or another group characteristic such as gender or ethnicity, conditional on student ability. If individuals with similar student ability experience similar causal effects of the subjective evaluation on outcomes of interest, then $\beta$ can also be interpreted as a measure of principal fairness, as recently introduced by \cite{imai2023principal}.

\subsection{Classical versus non-classical measurement error}

We cannot estimate \eqref{eqn:long} as we do not observe ability. Instead, students produce an unbiased but noisy signal of ability by taking a test, denoted by $\widetilde{s}_{t}$. The signal equals ability plus measurement error, $\widetilde{s_t}=s_t+m_t$, where $m_t$ is the signal's measurement error. An unbiased signal, meaning one that equals ability in expectation, implies that the error is classical.
\begin{assumption}\label{ass:classical}
Classical measurement error in the signal
\renewcommand{\theenumi}{\alph{enumi}}
\begin{enumerate}
    \item[] (Zero mean and cov) $\mathbb{E}[m_{t}]=0$ and $\mathrm{Cov}[m_{t},x]=0$ $\forall x \neq \widetilde{s}_t, t$. 
\end{enumerate} 
\end{assumption}
This classical error reflects, for instance, that students may guess answers, are more or less familiar with the specific selection of test items, or feel (un)well on the day of the test. According to classical testing theory, measures such as the fraction of correctly answered items follow \Cref{ass:classical} and are unbiased but noisy estimates of ability.

In modern assessment systems, test scores are rarely that simple. Instead, they are often the product of statistical IRT models designed to extract more information from the pattern of responses. For instance, \cite{jacob2016measurement} discuss that IRT test scores in many assessment systems in the US, such as the ECLS, the NELS:88, the ELS, and the HSLS are constructed from posterior means. Let $s^m_t$ denote the test score, which is then computed by 
\begin{align} \label{eqn:posterior}
   s_{t}^m=\mathbb{E}[s_t|\widetilde{s}_t] &=\int s_t f(s_t \mid \widetilde{s}_t) d s_t \\ \label{eqn:shrinking} 
   &\simeq  (1-\widetilde{\lambda} )\mathbb{E}[\widetilde{s}_t]+\widetilde{\lambda} \widetilde{s}_t, \quad \forall t.
\end{align} 
Students are treated as a random sample from a population of ability values $s_t$, characterized by the posterior density function $f(s_t \mid \widetilde{s}_t)$ from the IRT model. The posterior mean score is the mean of that population. Hence, instead of using the signal as the test score directly, the signal is used as a predictor for ability to generate a test score. This prediction is not perfect, so we can similarly write that the test score equals ability plus the score's measurement error, $s^m_{t}=s_{t}+{m}^{m}_{t}$. 

Posterior mean scores can be interpreted as Empirical Bayes estimates that, roughly, shrink the student’s (unbiased) maximum likelihood score toward the population mean in proportion to the noisiness of the maximum likelihood score \citep{jacob2016measurement}. Building on this interpretation, we approximate the posterior mean in \eqref{eqn:posterior} by the \cite{kelley1947fundamentals} estimator in \eqref{eqn:shrinking}, where $\widetilde{\lambda}$ is the shrinkage factor. With $\widetilde{\lambda}  \in [0,1)$, students with low (high) ability signal are pushed up (down) towards the population mean. As a result, the measurement error of the test score $m_t^m$ correlates negatively with ability. This feature of the non-classical error in \eqref{eqn:shrinking} is similar to the error generated by any posterior mean in \eqref{eqn:posterior}.

We will analyze contamination bias in $\beta$ when replacing ability in \eqref{eqn:long} by test scores measured through \eqref{eqn:shrinking}. Our results hold (approximately) for two types of test scores. The first are test scores estimated classically or through an IRT model using (weighted) maximum likelihood under a “fixed effect” approach. Both provide an unbiased estimate of ability. This corresponds to \eqref{eqn:shrinking} with $\widetilde{\lambda}=1$: the test score is equal to the signal, $s^m_t=\widetilde{s}_t$, and measurement error is classical, $m^{m}_t=m_t$.  The second are test scores estimated from an IRT model under a “random-effects” approach, where the scores are posterior means. Equation \eqref{eqn:shrinking} holds exactly only for the Kelley estimator, which corresponds to normality for both the prior and likelihood distributions. Generally, it is an approximation, and depending on the choice of these two distributions, the posterior mean in \eqref{eqn:posterior} may even lack a closed form. Because the Kelley estimator corresponds to an OLS regression of $s_t$ on $\widetilde{s}_t$, by the properties of OLS it provides the best linear approximation to the (possibly nonlinear) posterior mean. 

Our results speak less to settings where test scores are reported as posterior medians or modes, where multiple plausible values are drawn from a student’s posterior distribution, or where the prior distribution on ability incorporates student characteristics. The latter approach is used in the major US and international assessments NAEP and PISA. An alternative strategy is to jointly estimate the research model in \eqref{eqn:long} alongside the IRT model, treating the latter as a direct model of measurement \citep{junker2012use}. This requires access to item-level response data, which is typically unavailable. IRT models often produce noisier test scores at the tails of the distribution, reflecting the limited information available to these parametric methods when student ability does not align well with item difficulty. Our approach allows for the potential difference between so-called local and global reliability.

Non-classical measurement generated by \eqref{eqn:shrinking} with $\widetilde{\lambda} \in [0,1)$ is also considered by \cite{hyslop2001bias} in survey data. They consider a setting in which respondents are asked questions like ``what is the value of $s_t$?'' given an information set that consists of the unbiased but noisy signal $\widetilde{s}_t$. Respondents may passively report the ``flawed'' but unbiased value ($\widetilde{s}_t$), or may actively seek to provide an ``optimal'' response ($s^m_t$) given their information set. An argument for the latter is that respondents are likely aware of the lack of precision in the signal. Our results on the contamination bias below may thus also be relevant for the analysis of survey data more generally.

\subsection{Contamination bias}

\noindent The medium regression equation replaces $s_t$ with $s^m_t$,
\begin{align} \label{eqn:med}
    y_{t}= \beta^m SES + \gamma^m s_{t}^m+e_{t}^m.
\end{align}
To analyze what $\beta^m$ and $\gamma^m$ identify, we also introduce the balancing regression equation,
\begin{align}\label{eqn:balancing}
 \begin{pmatrix} s_t \\ s_{t}^m  \end{pmatrix}= \begin{pmatrix} s_t \\ \widetilde{\lambda} (s_t+m_t)  \end{pmatrix}=\begin{pmatrix} \delta \\ \delta^{m}  \end{pmatrix} SES +  \begin{pmatrix} u_{t} \\ u^m_{t} \end{pmatrix}.
\end{align} 
Note that $(1-\widetilde{\lambda} )\mathbb{E}[\widetilde{s}_t]$ is fixed and gets soaked up by the intercept of the regression. The following proposition formalizes what OLS identifies.
\begin{proposition}\label{prop:ols}
Under \Cref{ass:classical} and \Cref{eqn:shrinking}, it holds that: 
    \begin{align} 
    \label{eqn:deltam_ass}
    \delta^m=&\delta \widetilde{\lambda}, \\\label{eqn:gammam_ass}
    \gamma^m=&\gamma \Bigg(\frac{\lambda}{\widetilde{\lambda}}\Bigg), \\  
    \label{eqn:betam_ass}
    \beta^{m}=&\beta+\gamma\delta\big(1-\lambda\big).
    \end{align}    
With $\lambda=\Bigg(\frac{\sigma_{u_{t}}^2}{\sigma_{u_{t}}^2+\sigma_{m_{t}}^2} \Bigg)$, where $\sigma_{\bullet}^2$ denotes the variance of the variable in the subscript.
\end{proposition}

All proofs are provided in Appendix A to C. First consider \eqref{eqn:deltam_ass} and \eqref{eqn:gammam_ass} under classical measurement error, where $\widetilde{\lambda}=1$. Classical error on the left-hand side does not introduce bias in the estimate of $\delta^m$. However, as shown in \eqref{eqn:gammam_ass}, the estimate of $\gamma^m$ is subject to the well-known attenuation bias. The multivariate reliability (or signal-to-total variance) ratio of the test, denoted by $\lambda$, decreases from one toward zero as the measurement error variance increases, leading to attenuation bias in $\gamma^m$. Now consider the case of non-classical measurement error, where $\widetilde{\lambda} \in [0,1)$. \Cref{eqn:gammam_ass} shows that the attenuation bias in $\gamma^m$ may be smaller and disappears completely when the shrinkage factor equals the reliability ratio, $\widetilde{\lambda} = \lambda$. However, \eqref{eqn:deltam_ass} shows that non-classical error on the left-hand side results in attenuation of $\delta^m$.

\Cref{eqn:betam_ass} clarifies that both types of measurement error contaminate the estimate of $\beta^m$. Under classical measurement error, attenuation of $\gamma^m$ spills over to $\beta^m$. Under non-classical measurement error, the attenuation bias in $\gamma^m$ may be corrected, but at the expense of attenuation in $\delta^m$, which again spills over to $\beta^m$. Importantly, the contamination is identical for both types of measurement error, as the bias term involves the product of $\gamma^m$ and $\delta^m$, where the former is divided and the latter is multiplied by $\widetilde{\lambda}$. Although this may seem intuitive, to our knowledge this result has not been reported in the literature. In particular, \citet{hyslop2001bias} analyze \eqref{eqn:deltam_ass} and \eqref{eqn:gammam_ass} under both classical and non-classical measurement error, while \citet{pei2019poorly,gillen2019experimenting,ferman2022assessing} examine \eqref{eqn:betam_ass} in the case of classical measurement error only. Note that the contamination bias grows larger as $\gamma$ and $\delta$ increase and $\lambda$ decreases.
 
\subsection{The IV strategy}

The address measurement error, the instrumental variables (IV) strategy aims to use a second objective measure of ability in period $t$ as an instrument for the first measure. In practice, a lagged test score from period $t-1$ is often used instead \citep{botelho2015racial,ferman2022assessing,zhu2024}. Similar to above, students produce an unbiased but noisy signal of ability by taking a test in $t-1$, where $\widetilde{s}_{t-1} = s_{t-1}+m_{t-1}$. The corresponding test score satisfies \eqref{eqn:shrinking}, where $s_{t-1}^m=\mathbb{E}[s_{t-1}|\widetilde{s}_{t-1}]$ and $s^m_{t-1}=s_{t-1}+m^m_{t-1}$. We do not index the shrinking parameter $\widetilde{\lambda}$ by time, thereby assuming it remains constant across tests. Our main results hold without this simplification, though at the expense of additional notational clutter and reduced intuition. Using $s_{t-1}^m$ as an instrument for $s_{t}^m$ generates the following first and second stage regression equations respectively,
\begin{align}
    \label{eqn:fs}
    s_{t}^m &= \kappa SES + \pi s_{t-1}^m+e_t^{fs},\\
    \label{eqn:ss}
    y_t &= \beta^{iv} SES + \gamma^{iv} \widehat{s}_{t}^m +e_{t}^{iv}.
\end{align}

To show how this IV strategy can address measurement error, it is useful to also introduce the balancing regression for $t-1$,
\begin{align}\label{eqn:balancing_tmin1}
 \begin{pmatrix} s_{t-1} \\ s_{t-1}^m \end{pmatrix}= \begin{pmatrix} s_{t-1} \\ \widetilde{\lambda} (s_{t-1}+m_{t-1}) \end{pmatrix}=\begin{pmatrix} \delta_{t-1} \\ \delta_{t-1}^m   \end{pmatrix} SES +  \begin{pmatrix}  u_{t-1}  \\ u^m_{t-1} \end{pmatrix},
\end{align} 
and the variable $\Delta_{s_{t}}=s_t-s_{t-1}$, which captures the change in unobserved ability between time period $t$ and $t-1$. We describe two additional assumptions that place increasing restrictions on $\Delta_{s_{t}}$. 
\begin{assumption}\label{ass:fs}
The covariance between $\Delta_{s_t}$, $SES$ and $s_{t}$ is zero 
\renewcommand{\theenumi}{\alph{enumi}}
\begin{enumerate}
    \item (SES) $\mathrm{Cov}[\Delta{s_{t}},SES]=0$.
    \item (Ability) $\mathrm{Cov}[\Delta_{s_{t}},s_{t}]=0$.
    \end{enumerate} 
\end{assumption}
\noindent \Cref{ass:fs}(a) requires the change in ability to be the same, on average, across low and high SES students. \Cref{ass:fs}(b) requires this for low and high ability students in time $t$. 
\begin{assumption}\label{ass:rf}
Constant $\Delta_{s_t}$
\renewcommand{\theenumi}{\alph{enumi}}
\begin{enumerate}
    \item[] (Constant) $\Delta{s_{t}}$ is the same for each student.
    \end{enumerate} 
\end{assumption}
\noindent \Cref{ass:rf} requires that the change in ability is the same for each student, which is stronger than \Cref{ass:fs}. The following proposition summarizes what the IV strategy identifies while combining \Cref{ass:classical} with either \ref{ass:fs} or \ref{ass:rf}.
\begin{proposition}\label{prop:iv}
\indent (i) Under \Cref{ass:classical} and \ref{ass:fs}, and \Cref{eqn:shrinking}, it holds that: 
    \begin{align}
       \pi'=&\lambda, \\
         \gamma^{iv}=&\gamma \Bigg(\frac{1}{\widetilde{\lambda}}\Bigg) B_{\gamma}, \\
        \beta^{iv}=& \beta+\gamma\delta\bigg(1-B_{\gamma}\bigg).    
    \end{align} 
With $\pi'=\pi \Bigg( \frac{\sigma_{u_{t}}^2+ \sigma_{\Delta_{s_t}}^2+\sigma_{m_{t-1}}^2}{\sigma_{u_{t}}^2+\sigma_{m_{t}}^2}\Bigg)$ and $B_{\gamma}= 1 -  \Bigg(\frac{\mathrm{Cov}[e_{t},\Delta_{s_t}]}{\mathrm{Cov}[y_{t},u_{t}]}\Bigg)$.  

\vspace*{.2in}
\indent (ii) Under \Cref{ass:classical} and \ref{ass:rf}, and \Cref{eqn:shrinking}, it holds that: 
    \begin{align}
        \pi'=&\lambda, \\
        \gamma^{iv}=&\gamma \Bigg(\frac{1}{\widetilde{\lambda}}\Bigg), \\
        \label{eqn:betaiv_ass3}
        \beta^{iv}=&\beta.          
    \end{align}
With $\pi'=\pi \Bigg( \frac{\sigma_{u_{t}}^2+\sigma_{m_{t-1}}^2}{\sigma_{u_{t}}^2+\sigma_{m_{t}}^2}\Bigg)$.  
\end{proposition}

Under classical measurement error with $\widetilde{\lambda}=1$, \Cref{prop:iv} formalizes that the IV strategy yields consistent estimates under \Cref{ass:rf}. This result is not new. For instance, \cite{gillen2019experimenting} use IV to address contamination bias from classical error in \eqref{eqn:betaiv_ass3} and review similar applications in the literature. Intuitively, a constant $\Delta{s_{t}}$ implies that the test score in period $t-1$ is truly a second objective measure of ability in period $t$. The first stage regression of $s_{t}^m$ upon $s_{t-1}^m$ measures to which extent this relationship is diluted due to measurement error, so that it reveals the reliability ratio $\lambda$.
%
Furthermore, the reduced form equation is similar to the medium regression equation \eqref{eqn:med} while replacing $s_{t}^m$ with $s_{t-1}^m$. With a constant $\Delta{s_{t}}$ the reduced form coefficient reveals the estimate $\gamma^m=\gamma\lambda$, which was attenuated towards zero by the reliability ratio. Dividing the reduced form by the first stage corrects for this attenuation. Since $\gamma^{iv}=\gamma$ and $\delta^m=\delta$, we also have that $\beta^{iv}=\beta$.

In the presence of non-classical error with $\widetilde{\lambda}\in[0,1)$, the IV strategy also addresses contamination bias under \Cref{ass:rf}, though for a different reason. To illustrate, suppose the shrinkage factor equals the reliability ratio, such that $\widetilde{\lambda}=\lambda$. In this case, the OLS coefficient on test scores $\gamma^m$ does not suffer from attenuation bias. However, the IV strategy still divides by the first stage, which means that the IV coefficient $\gamma^{iv}$ is overestimated by a factor equal to the inverse of the reliability ratio. As noted by \cite{hyslop2001bias}, IV has a bias away from zero in the presence of non-classical measurement error.
%
However, this upward bias in $\gamma^{iv}$ is exactly what is needed to correct for the contamination bias to $\beta^{iv}$, which arises from attenuation in the estimate of $\delta^m$ from the balancing regression. Since the contamination bias term involves the product of the test score and balancing coefficients, the upward bias in the former offsets the downward bias in the latter. As a result, the IV strategy also effectively addresses the contamination bias in the presence of non-classical measurement error.

\subsection{The EIV strategy}

Based on the results of \Cref{prop:ols}, the errors-in-variables (EIV) strategy directly formulates expressions for the parameters of interest,
\begin{align}\label{eqn:gamma_est}
    \gamma^{eiv}=&\frac{\gamma^m}{\lambda}, \\  
    \label{eqn:beta_est}
    \beta^{eiv}=&\beta^m-\gamma^m \delta^m \Bigg(\frac{1-\lambda}{\lambda}\Bigg).
\end{align}
All the elements on the right-hand side are known under \Cref{ass:classical}, except for the reliability ratio. 
The EIV strategy then plugs a reliable estimate for $\lambda$ into \eqref{eqn:gamma_est} and \eqref{eqn:beta_est} to estimate the parameters of interest. 

\subsubsection{The first stage}

\Cref{prop:iv} shows that the first stage can be used as an estimate for $\lambda$ under \Cref{ass:fs}. The intuition for this result follows more easily in a  univariate setting without $SES$ as an independent variable. The univariate reliability ratio is defined by replacing $\sigma_{u_{t}}^2$ in $\lambda$ with $\sigma_{s_{t}}^2$. \Cref{ass:fs}(b) implies that $\mathrm{Cov}[s_t,s_{t-1}]=\sigma_{s_{t}}^2$, so the coefficient from the univariate first stage regression of $s^m_{t}$ on $s^m_{t-1}$ is proportional to the univariate reliability ratio. This result extends to the multivariate setting since together \Cref{ass:fs}(a)-(b) imply that $\mathrm{Cov}[u_t,u_{t-1}]=\sigma_{u_{t}}^2$.

Our EIV first stage (EIV FS) strategy uses the first stage as an estimate for $\lambda$. The following proposition formalizes what EIV FS identifies under \Cref{ass:classical} and \ref{ass:fs}.
\begin{proposition}\label{prop:fs}
Under \Cref{ass:classical} and \ref{ass:fs}, and \Cref{eqn:shrinking}, it holds that: 
    \begin{align}
        \gamma^{eiv}=&\gamma \Bigg(\frac{1}{\widetilde{\lambda}}\Bigg), \\ 
        \beta^{eiv}=&\beta.
    \end{align}
Where $\gamma^{eiv}$ and $\beta^{eiv}$ are defined by replacing $\lambda$ with $\pi'$ in \eqref{eqn:gamma_est} and \eqref{eqn:beta_est}.    
\end{proposition}

Similar to the IV strategy, EIV FS corrects attenuation bias in the test score coefficient under classical measurement error, overestimates it under non-classical error, and addresses contamination bias in the SES coefficient in both cases. Unlike the IV strategy, EIV FS achieves these results under the weaker \Cref{ass:fs}. It relies directly on the medium regression estimate, whereas the IV strategy uses the reduced form estimate. The two are equivalent only if $\mathrm{Cov}[e_{t},\Delta_{s_t}]=0$, which 
%$\mathrm{Cov}[e_{t},\Delta_{s_t}]=-\mathrm{Cov}[e_{t},s_{t-1}]=0$
holds when $\Delta_{s_t}$ is constant. Whereas \Cref{ass:rf} only seems likely when the two tests are administered arbitrarily shortly after one another, \Cref{ass:fs} appears generally more plausible. In particular, under \Cref{ass:fs} the teachers may pay attention to $\Delta_{s_t}$ in their subjective evaluations, so that it is part of the error term $e_{t}$, but $\Delta_{s_t}$ cannot correlate with the included variable $SES$ and unobserved ability $s_t$. %\Cref{ass:fs} does not require that $s_{t-1}$ is linearly related to $s_t$ as in \cite{ferman2022assessing}.

\subsubsection{Test score history}

We propose a novel EIV test score history (EIV TH) strategy that recovers the reliability ratio under \Cref{ass:classical} only, while making use of data on students' test score histories. The intuition of this strategy is to approximate the ideal conditions of the test-retest method, which estimates the reliability ratio using a test administered at time $t-\epsilon$, arbitrarily short before time $t$ \citep{Cito2017}. With this ideal test from period $t-\epsilon$, the first stage equation is,
\begin{align}
    s_{t}^m=\kappa_{t-\epsilon}SES +\pi_{t-\epsilon} s_{t-\epsilon}^m+e_t^{fs_{\epsilon}},
\end{align}
where we include a time subscript on the coefficients. For a test administered arbitrarily close to time $t$, \Cref{ass:rf} holds trivially, and \Cref{prop:iv} shows that $\pi_{t-\epsilon}'= \lambda$. While a test score that close to time $t$  may not be available, we may observe multiple test scores further away. It is plausible that such scores satisfy \Cref{ass:classical}, even if \Cref{ass:fs} and \ref{ass:rf} may not hold. Nevertheless, the EIV TH strategy leverages these test scores observed at earlier points in time to identify the reliability ratio.

We begin by repeating the first stage equation, sequentially replacing the test score from period $t-1$ on the right-hand side with scores from earlier periods ,
\begin{align}\label{eqn:fs_tau}
    s_{t}^m= \kappa_{t-\tau}SES + \pi_{t-\tau} s_{t-\tau}^m+e_t^{fs_{\tau}}, \quad \forall \tau \geq 1.
\end{align}
It follows from  the proof of \Cref{prop:iv} that under \Cref{ass:classical} only, $\pi_{t-\tau}'$ is a biased estimate for the reliability ratio,
\begin{align}
   \pi_{t-\tau}'=& \lambda B_{\pi_{t-\tau}}, \quad \forall \tau \geq 1,
\end{align}
with $B_{\pi_{t-\tau}}=\Bigg( \frac{\mathrm{Cov}[u_{t},u_{t-\tau}]}{\sigma_{u_{t}}^2}\Bigg)$. Therefore, we use $\pi_{t-\tau}'$ as a dependent variable in a forecasting regression equation of the form,
\begin{align}\label{eqn:forecasting}
    \pi_{t-\tau}'=& f(t-\tau)+\varepsilon_{t-\tau}. 
\end{align}
Subsequently, we make a one-period out-of-sample forecast towards $\tau=\epsilon$, denoted by $ \widehat{f}(t-\epsilon)$. Without forecasting error this resembles the ideal conditions of the test-retest method, such that $\widehat{f}(t-\epsilon)={f}(t-\epsilon)= \pi_{t-\epsilon}'=\lambda$. 
In the absence of forecasting error, the chosen polynomial in \eqref{eqn:forecasting} accurately captures the changes in $\pi'_{t-\tau}$ over time. While the EIV FS strategy estimates $\gamma^{eiv}$ and $\beta^{eiv}$ by replacing $\lambda$ with $\pi'$ in \eqref{eqn:gamma_est} and \eqref{eqn:beta_est}, EIV TH estimates them by replacing $\lambda$ with ${f}(t-\epsilon)$.

\begin{comment}

Cameron and Trivedi (textbook, 26.3) discuss both IV and errors-in-variables model as a solution to measurement error, but they do not make the connection between the two. They discuss IV in the general case: you need a Z that correlates with D but not with epsilon. They discuss the errors-in-variables model as something that they call "data replication": you need two measurements of the same underlying concept to estimate the signal to noise ratio, and then you can correct OLS manually, which is exactly what the Stata command eivreg does. The connection between the two is that the first stage in an IV method is meant to estimate the signal to noise ratio.

card (1996): interested in the effect of unions on wages and discusses that measurement error in union reporting can be solved  via a two-step estimation approach if you know some auxilary parameters, such as the union status missclafication rate. he does not focus on the contimation bias, but only on the main effect of unions. he does not make the link with IV, but estimates the auxiliary parameters with auxiliary data. 
Krashinsky (2004): in spirit of card, discusses that measurement error can be solved via a measurement-error-corrected GLS, CRE, or FE estimator (this is similar to eivreg) if you know the variance covariance matrix of the measurement errors. He only focuses on the main effect and does not make the link with IV. He obtains estimates for the variance covariance matrix from other sources/papers, where one souce is the "data replication" aproach.  
Ashenfelter and Krueger (1994): estimate the returns to schooling with twin fe and discuss measurement error. with a fe estimator this problem is excerbated. they solve this with an IV approach where the instrument is the reported years of schooling by the twin. they only focus on the main effect. they do not make the link that the first stage is all you need.

\end{comment}



